\section{总结与延伸}

\subsection{全讲框架回顾}

\begin{table}[H]
\centering
\begin{tabular}{p{2.6cm}p{4.2cm}p{4.1cm}}
\toprule
\textbf{模块} & \textbf{核心结论} & \textbf{实践动作} \\
\midrule
问题定义 & Agent 后训练目标是可验证任务完成 & 建立 preference + verifier 联合目标 \\
任务建模 & 核心三元组：环境、工具、验证器 & 数据与系统日志按轨迹组织 \\
数据工程 & 难点在组合覆盖而非样本总量 & 扩展环境/工具/验证组合，显式纳入失败样本 \\
评估体系 & 必须 Holistic，防止指标幻觉 & 多 verifier + harness swapping + 污染监控 \\
训练流程 & Light SFT 打底，RL 负责突破 & 控制 SFT 强度，保留探索空间 \\
RL 关键点 & 防熵塌缩、难度调度、多样高质轨迹 & 监控熵与难度分层，优化采样与反馈闭环 \\
算法实现 & on-policy、更新强度平衡、熵正则有用 & 保留消融与审计，避免技巧神化 \\
\bottomrule
\end{tabular}
\caption{Post-Training Verifiable Agents 的工程化总图}
\end{table}

\begin{importantbox}{一句话总结}
可验证 Agent 训练的本质，不是 ``让模型更会说''，而是 ``让模型在可观测、可验证、可治理的环境中持续学会做对事''。
\end{importantbox}

\subsection{延伸阅读与实践建议}

前面的总结表给出机制关系，本小节进一步把它们转成可复现实验。阅读论文时应同时复现环境、verifier、采样预算和基线，实践中则优先做 harness swap、SFT 强度消融、entropy 监控与轨迹级数据治理，避免只复制算法名称。

\begin{longtable}{p{3.0cm}p{8.0cm}}
\toprule
\textbf{主题} & \textbf{建议阅读/实践方向} \\
\midrule
\endfirsthead
\toprule
\textbf{主题} & \textbf{建议阅读/实践方向} \\
\midrule
\endhead
课程原视频 & Berkeley RDI 发布的本讲完整视频（Jiantao Jiao, Post-Training Verifiable Agents），建议结合本笔记的关键时间戳回看图示部分。 \\
评估工程 & 重点补齐 verifier 分层设计：结果层、结构层、过程层、安全层；在内部基准上进行 harness swapping 实验。 \\
训练策略 & 对比 ``重 SFT + 轻 RL'' 与 ``轻 SFT + 强 RL'' 两种配置，追踪熵变化、成功率和 OOD 泛化差异。 \\
熵与探索 & 针对熵塌缩建立日常监控面板：token entropy、轨迹去重率、有效探索比例、失败类型分布。 \\
数据体系 & 从 ``按任务存样本'' 升级到 ``按轨迹存状态转移''，将工具调用参数与中间反馈纳入训练可回放数据。 \\
治理闭环 & 引入权限边界、调用审计、异常回滚机制，把部署可控性作为后训练优化目标的一部分。 \\
\bottomrule
\end{longtable}

\subsection{进一步思考}
\begin{itemize}
    \item 当 verifier 无法覆盖复杂开放任务时，如何避免模型在弱约束区域形成投机策略？
    \item 在成本受限条件下，如何选择最具信息增益的轨迹进行 on-policy 更新？
    \item 可验证 Agent 的评估是否应从 ``单模型分数'' 转向 ``系统级可靠性曲线''？
    \item 如果训练目标同时包含能力与治理，损失函数和评估基准应如何共同设计？
\end{itemize}

%% --- 正文结束 --- %%

\end{document}
